\noindent326\hfill\textbf{CORRESPONDENCE.\quad ON NETWORK ADJUSTMENT.}
\begin{center}\textsc{Gauss} to \textsc{Gerling}.\quad Göttingen, 14 November 1838.\end{center}

.\ .\ .\ .\ .\ .\ As for the formulas for the number of condition equations,
I have no objection at all to their publication. But it is impossible
for me \textls[120]{now} to specify the exceptional cases precisely and
exhaustively, and I do not know whether that will be \textls[120]{possible}
at all without going into the matter more deeply. S[ine] m[editatione]
it should suffice to say that if in a system of $p$ points joined by $l$
connecting lines the direction of each connecting line is measured at
\textls[120]{both} endpoints, that is, tied to another direction line of
the system, there must be

$l-p+1$ mutually independent condition equations from angle sums alone
and, in addition,

$l-2p+3$ \textls[120]{other} condition equations, which in the \textls[120]{most usual}
cases all fall into the category of those given in
Supplem. Theor.\ p.\ 31 under III[\textsuperscript{*})].
\par\medskip
\noindent\rule{.24\textwidth}{.4pt}\par\smallskip
{\footnotesize\noindent\hspace*{1.5em}[\textsuperscript{*}) Volume IV, p.\ 81.]\par}
\vspace{2em}
\begin{center}\rule{.18\textwidth}{.4pt}\end{center}
